%==============================================================================
\section{Model}
\label{sec:model}
%==============================================================================

The model has four stages that capture the blockholder's information advantage, the market's inference problem, and the bidder's response.

%------------------------------------------------------------------------------
\subsection{Timeline}
%------------------------------------------------------------------------------

Time is discrete, $t \in \{0, 1, 1.5, 2\}$.

\begin{enumerate}[leftmargin=2em]
\item[$t=0$:] Nature draws a standalone fundamental $v$. The blockholder observes a private signal $s$ about $v$.
\item[$t=1$:] The blockholder chooses an action $(q,a)$ from the feasible set $\{(-1,0),(0,0),(0,1),(+1,1)\}$, where $q \in \{-1,0,+1\}$ is the trade and $a \in \{0,1\}$ is engagement. A noise trader submits $z \in \{-1, 0, +1\}$. The market maker observes total order flow $X = q + z$ and disclosure indicator $D$, then sets a competitive price $P(X, D)$.
\item[$t=1.5$:] A potential bidder observes $(P(X,D), D)$ and a private synergy shock $\xi$, then decides whether to initiate a takeover attempt.
\item[$t=2$:] Payoffs are realized: either the takeover is consummated at the offered price, or the firm remains standalone with (possibly) improved fundamentals due to engagement.
\end{enumerate}

%------------------------------------------------------------------------------
\subsection{Notation Summary}
%------------------------------------------------------------------------------

Table~\ref{tab:notation} in Appendix~\ref{app:notation} provides a quick reference for notation. The remainder of this section introduces each model component in turn, beginning with the blockholder's information environment and ending with the equilibrium concept.

%------------------------------------------------------------------------------
\subsection{Fundamentals and Information}
%------------------------------------------------------------------------------

The standalone fundamental is $v \sim \mathcal{N}(\mu, \sigma_v^2)$. The blockholder observes a private signal
\[
s = v + \varepsilon, \quad \varepsilon \sim \mathcal{N}(0, \sigma_\varepsilon^2),
\]
independent of $v$. Let $\sigma_s^2 \equiv \sigma_v^2 + \sigma_\varepsilon^2$. The blockholder's posterior mean is linear:
\[
\hat{v}(s) \equiv \E[v \mid s] = \mu + \beta(s - \mu), \quad \beta \equiv \frac{\sigma_v^2}{\sigma_v^2 + \sigma_\varepsilon^2} \in (0,1).
\]

The bidder draws an independent synergy shock $\xi \sim \mathcal{N}(0, \sigma_\xi^2)$. This signal structure gives the blockholder a transient information advantage that she can exploit through trading, engagement, or both.

%------------------------------------------------------------------------------
\subsection{Trading and Liquidity}
%------------------------------------------------------------------------------

The blockholder has an initial stake of one share (normalization). Her order $q \in \{-1, 0, +1\}$ implies an end-of-$t=1$ holding of $h = 1 + q$ shares.

Noise trading is discrete:
\[
z \in \{-1, 0, 1\}, \quad \PP(z = 0) = 1 - \kappa, \quad \PP(z = \pm 1) = \kappa/2,
\]
where $\kappa \in (0, 1)$ indexes noise-trading intensity (market liquidity). Higher $\kappa$ means more frequent nonzero noise imbalance and therefore weaker inference from order flow.

Total order flow observed by the market maker is
\[
X = q + z \in \{-2, -1, 0, 1, 2\}.
\]

The blockholder's trading decision is intertwined with her engagement choice, to which I now turn.

%------------------------------------------------------------------------------
\subsection{Disclosure Rule}
%------------------------------------------------------------------------------

Let $D \in \{0, 1\}$ be a public indicator for whether a stake acquisition crosses a disclosure threshold. I model disclosure as \emph{stake-triggered}: in the discrete framework,
\[
D = 1 \quad \Longleftrightarrow \quad (q = +1).
\]
That is, \textbf{if the blockholder acquires an additional share, a threshold-crossing disclosure occurs}. This captures the institutional logic that crossing an ownership threshold triggers a regulatory filing (e.g., Schedule 13D), independent of whether the filing explicitly reflects intent. Below-threshold engagement remains hidden.

\paragraph{Timing of disclosure.} The disclosure indicator $D$ is determined simultaneously with the blockholder's action $(q, a)$ and is observed by the market maker at the same time as order flow $X$. Specifically, the sequence within $t=1$ is:
\begin{enumerate}[label=(\roman*),leftmargin=2em,itemsep=0pt]
\item The blockholder chooses $(q, a)$, which determines $D = \1\{q=+1\}$.
\item Noise trader submits $z$; total order flow $X = q + z$ is realized.
\item The market maker observes $(X, D)$ jointly and sets price $P(X, D)$.
\end{enumerate}
This timing ensures that disclosure does not reveal additional information beyond what is captured by $(X, D)$ jointly.

%------------------------------------------------------------------------------
\subsection{Engagement Technology (Voice)}
\label{sec:engagement}
%------------------------------------------------------------------------------

If the blockholder chooses engagement $a = 1$, she pays a private, signal-dependent cost $C(s) > 0$ at $t = 1$. This cost is observed by the blockholder when she observes $s$ but not by the market maker or bidder. I assume engagement costs are decreasing in the signal---a blockholder with more favorable private information finds it easier to build a compelling case for change. I parameterize the cost as
\[
C(s) = C_0 \exp\!\left(-\chi \frac{s - \mu}{\sigma_s}\right),
\]
where $C_0 > 0$ is the cost at the prior mean and $\chi \geq 0$ controls cost sensitivity to signal quality (Assumption~A2). This monotone, strictly positive specification ensures that neither engagement nor passivity is always optimal when the blockholder holds her initial stake, yielding an interior cutoff structure (Assumption~A1).%
\footnote{Sufficient conditions for interior cutoffs include $C(\mu) < \delta(\tilde{\Delta} + (\tilde{m}-m_0))$ (engagement is worthwhile at the prior mean) and $\lim_{s\to-\infty} C(s) > \delta(\tilde{\Delta} + (\tilde{m}-m_0))$ (engagement is dominated for sufficiently low signals), together with the monotonicity of $C(s)$. I also assume that stake-building (Public Voice) is not always dominated when engagement occurs, so that disclosure is triggered for sufficiently high signals. When $\chi > 0$, cost heterogeneity supports an interior Passive Hold region; when $\chi = 0$, costs are constant and the equilibrium collapses to the two-cutoff case.}

Engagement succeeds with probability $\rho \in (0,1]$ and affects outcomes in two ways when successful:

\paragraph{Outside option improvement.} If engagement succeeds and no takeover is consummated, the standalone payoff is increased by $\Delta > 0$ at $t=2$. If engagement fails, standalone payoff remains $v$.

\paragraph{Bargaining/resistance premium.} If engagement succeeds and a takeover is consummated, the required premium increases from $m_0$ to $m_1$, where $m_1 > m_0 \geq 0$. The wedge captures bargaining power, defensive tactics, or any friction that raises the price needed to complete the acquisition when an engaged blockholder is present. Failed engagement yields the baseline premium $m_0$.
I interpret $m_0$ and $m_1$ as \emph{per-share takeover premia} above the market price (so the consummated offer satisfies $b=P+m$); they are distinct from $\hat{v}(s)$, the posterior mean of fundamentals.

\paragraph{Risk-neutral simplification.} Since success is unobserved before $t=2$ and all parties are risk-neutral, I define expected engagement effects:
\[
\tilde{\Delta} \equiv \rho \Delta, \qquad \tilde{m} \equiv m_0 + \rho(m_1 - m_0).
\]
Then engagement ($a=1$) yields expected standalone improvement $\tilde{\Delta}$ and expected premium wedge $\tilde{m}$. I assume the premium wedge satisfies $\tilde{m} > m_0$, which requires $\rho(m_1 - m_0) > 0$ (Assumption~A3). This ensures that engagement has real consequences for takeover outcomes---an engaged blockholder can extract higher premia through an improved bargaining position or defensive tactics.

%------------------------------------------------------------------------------
\subsection{Bidder Entry and Bidding Terms}
%------------------------------------------------------------------------------

The bidder follows a price-contingent entry rule as in standard takeover models \citep{GrossmanHart1980,Fishman1988,BulowHuangKlemperer1999,EdmansGoldsteinJiang2012}. The bidder observes the public pair $(P(X,D), D)$ and a private synergy shock $\xi$. On the non-disclosed branch $D=0$, I assume the equilibrium price schedule separates order-flow realizations---$P^*(X,0) \neq P^*(X',0)$ for all distinct $X, X' \in \{-2,-1,0,1\}$---so that the bidder can infer $X$ from $(P,0)$ (Assumption~A7).%
\footnote{Injectivity on the $D=0$ branch is a standard regularity condition in discrete-order-flow models: different $X$ imply different posteriors about fundamentals and hence different competitive prices. Assumption~A7 avoids a notational distinction between conditioning on $(P,D)$ and conditioning on $(X,D)$.}
I therefore use $p(X,0)$ and $m(X,0)$ as shorthand for $p(P(X,0),0)$ and $m(P(X,0),0)$. On the disclosed branch $D=1$, disclosure pins down Public Voice and hence $a=1$, so $X$ reveals only the noise realization and carries no information about fundamentals. Under Assumption~A5 (described below), this implies that equilibrium beliefs and the pricing fixed point are identical across $X\in\{0,1,2\}$ on the disclosed branch, so $P^*(X,1)$ and $p(X,1)$ are constant across $X$ (see Appendix~\ref{app:proof-disclosed-invariance}). I therefore continue to write $p(X,1)$ as shorthand for $p(P^*(X,1),1)$.

I assume that synergies and costs satisfy $\bar{S}-K > m_0-\mu$ and $\bar{S}-K < m_1+(\mu+3\sigma_v)$ (Assumption~A4). These bounds ensure that takeover bids occur with interior probability: the bidder is neither always deterred nor always willing to bid regardless of prices.

Conditional on $(P,D)$, the bidder's net deal surplus is
\[
\Pi_B(X,D) = \bar{S} - P(X,D) + \xi - m(X,D) - K,
\]
where $\bar{S}$ is baseline synergy, $K > 0$ is a fixed bidding cost, and $m(X,D)$ is the expected premium wedge conditional on $(X,D)$. The market price enters surplus dollar-for-dollar because the bidder must pay the target's market value (plus the premium wedge) to complete the acquisition.

The premium wedge ultimately depends on whether the blockholder has engaged. I define the premium required to complete the acquisition as a function of the \emph{realized} engagement action $a$:
\[
m(a) \equiv m_0 + a\cdot(\tilde{m}-m_0), \qquad a\in\{0,1\}.
\]
I interpret $m(a)$ as the (ex-ante) premium required at closing, determined by the blockholder's engagement-driven resistance.
When engagement is not publicly disclosed, the market maker and bidder do not observe $a$ and instead form the posterior $\pi(X,D)=\PP(a=1 \mid X,D)$. The corresponding \emph{expected} premium wedge conditional on $(X,D)$ is therefore
\[
m(X,D) \equiv \E[m(a) \mid X,D] = m_0 + (\tilde{m} - m_0)\cdot \pi(X,D), \quad \pi(X,D) \equiv \PP(a = 1 \mid X, D).
\]

I define $m(P,D)\equiv m(X(P,D),D)$ and the takeover completion probability conditional on observed $(P,D)$ as
\begin{equation}
p(P,D) \equiv 1 - \Phi\left(\frac{m(P,D) + K - \bar{S} + P}{\sigma_\xi}\right),
\label{eq:bid-prob}
\end{equation}
and use the shorthand $p(X,D)\equiv p(P(X,D),D)$ in equilibrium. The bidder initiates a takeover attempt if and only if $\Pi_B(X,D) \geq 0$, so $p(P,D)$ is the probability of a consummated transaction.
Here $\Phi(\cdot)$ is the standard normal cumulative distribution function.

If the takeover is consummated, the \emph{realized} per-share offer is
\[
b(X,D,a) = P(X,D) + m(a).
\]

%------------------------------------------------------------------------------
\subsection{Terminal Payoff and Price Formation}
%------------------------------------------------------------------------------

Let $\delta = \exp(-r\tau)$ (approximately $1/(1+r)$ for short horizons) be the discount factor between the pricing/entry date and payoff/settlement at $t=2$, where $\tau$ is the time between these dates \citep{Cochrane2005,Duffie2010}.%
\footnote{If one measures payoffs in time-$1$ present-value units, $\delta$ can be set to one and later reintroduced by scaling. I keep $\delta$ explicit throughout.}

If no takeover is consummated, per-share terminal payoff is $y^{\text{stand}} = v + a\tilde{\Delta}$. If a takeover is consummated, per-share payoff is $y^{\text{M\&A}} = b(X,D,a) = P(X,D) + m(a)$.

Therefore, the per-share terminal payoff is
\[
Y = \1\{\textup{bid}\} \cdot (P(X,D) + m(a)) + (1 - \1\{\textup{bid}\}) \cdot (v + a\tilde{\Delta}),
\]
where $\1\{\textup{bid}\}$ denotes a \emph{consummated} takeover rather than a mere offer. Bargaining, board resistance, and regulatory frictions are captured in reduced form by the premium wedge $m(\cdot)$ and the fixed cost $K$. A simple robustness extension would multiply $p(P,D)$ by a completion probability $\alpha(P,D)\in(0,1]$ in the pricing fixed point, without altering the main mechanism.

Conditional on $(X,D)$, $\E[m(a) \mid X,D]=m(X,D)$, so the expected takeover payoff equals $P(X,D)+m(X,D)$; this keeps the pricing fixed-point equations below unchanged.

The market maker is competitive and sets
\begin{equation}
P(X,D) = \delta \, \E[Y \mid X, D].
\label{eq:pricing}
\end{equation}

Equation~\eqref{eq:pricing} is the key equilibrium discipline: the price reflects the possibility of takeover and the premium wedge, and takeover decisions depend on that same price. I assume that for each information set $(X,D)$ and each feasible posterior $\pi(X,D)$, the competitive pricing fixed point admits a unique solution $P^*(X,D)$, and that takeover entry is not too sensitive to prices: $\delta\cdot\sup_P|\partial p/\partial P|<1$ (Assumption~A5).%
\footnote{Assumption~A5 bounds the strength of the price-to-entry feedback loop and ensures that comparative statics are well-defined. Under the Gaussian bidding rule, a sufficient condition is $\delta/\sigma_\xi < 1/\phi(0)$, where $\phi(\cdot)$ is the standard normal density.}

%------------------------------------------------------------------------------
\subsection{Blockholder Payoff and Equilibrium Concept}
%------------------------------------------------------------------------------

Conditional on signal $s$, the blockholder chooses $(q, a)$ to maximize expected present value. If she submits $q$ and ends with $h = 1 + q$ shares, her $t = 1$ value is
\[
U(q, a \mid s) = \E\Big[-q \cdot P(X,D) + \delta \cdot h \cdot Y - a \cdot C(s) \,\Big|\, s, q, a\Big],
\]
where $-qP(X,D)$ is the net cash flow from trading at $t = 1$ (selling $q = -1$ yields $+P$; buying $q = +1$ yields $-P$).

\begin{definition}[Perfect Bayesian Equilibrium]
A \emph{Perfect Bayesian Equilibrium} consists of:
\begin{enumerate}[label=(\roman{enumi})]
\item a blockholder strategy mapping $s$ into $(q, a)$,
\item beliefs $\pi(X,D) = \PP(a = 1 \mid X, D)$ consistent with Bayes' rule,
\item a bidder entry strategy satisfying~\eqref{eq:bid-prob}, and
\item a competitive price schedule $P(X,D)$ satisfying~\eqref{eq:pricing}.
\end{enumerate}
\end{definition}

Because noise trading $z$ has full support on $\{-1,0,+1\}$, every $(X,D)$ pair that is feasible under an on-path action occurs with positive probability. In particular, when $D=0$, $X\in\{-2,-1,0,1\}$; when $D=1$, only Public Voice is feasible so $X\in\{0,1,2\}$. Consequently, Bayes' rule pins down $\pi(X,D)$ for all reached information sets; off-path beliefs are irrelevant.%
\footnote{Throughout, I maintain Assumptions~A1--A7 (the Standing Assumptions). These ensure interior action regions, non-trivial takeover outcomes, and well-behaved equilibrium existence. I verify each assumption numerically in the baseline calibration. Assumption~A6 (contraction) imposes that, on a suitable compact subset of cutoff vectors, the equilibrium mapping is a contraction, guaranteeing existence and uniqueness of the monotone-cutoff equilibrium. Assumption~A8, which provides sufficient conditions for non-monotonicity of minority gains in liquidity, is stated where it is first invoked in the equilibrium analysis.}
